% --- Archon physics formalization source begin ---
% archon:physics
% archon:covers PhyXMiniProblems/problem_phyx_mini_0539.lean
% archon:source-report reports/phyx_mini/problem_phyx_mini_0539.source.json

\chapter{Physics problem phyx\_mini\_0539}
\label{ch:PhyXMiniProblems_problem_phyx_mini_0539}

\paragraph{Problem source.}
\#\# Physical scenario

Assume a photomultiplier tube for detecting radiation has seven dynodes with potentials of \textbackslash{}(100, 200, 300, \textbackslash{}ldots, 700 \textbackslash{}text\{ V\}\textbackslash{}) as shown in figure. The average energy required to free an electron from the dynode surface is \textbackslash{}(10.0 \textbackslash{}text\{ eV\}\textbackslash{}). Assume only one electron is incident and the tube functions with \textbackslash{}(100\textbackslash{}\%\textbackslash{}) efficiency.

\#\# Question

What is the energy available to the counter for all the electrons arriving at the last dynode?

\#\# Answer choices

- A: A: \textbackslash{}(10\textasciicircum{}10 \textbackslash{}text\{ eV\}\textbackslash{})

- B: B: \textbackslash{}(10\textasciicircum{}6 \textbackslash{}text\{ eV\}\textbackslash{})

- C: C: \textbackslash{}(10\textasciicircum{}12 \textbackslash{}text\{ eV\}\textbackslash{})

- D: D: \textbackslash{}(10\textasciicircum{}8 \textbackslash{}text\{ eV\}\textbackslash{})

\#\# Figure caption (auxiliary; use the image as primary evidence)

The image shows a diagram of a photomultiplier tube. Here is the detailed description:

- **Scintillation Crystal:** At the top, there is a scintillation crystal, which is the starting point for the incident particle or photon interaction, shown with a red streak.
  
- **Photocathode:** Below the scintillation crystal, a photocathode is illustrated. The interaction from the crystal leads to the production of electrons here.

- **Voltage Stages:** The main body of the tube is divided into several stages, each labeled with increasing voltage levels:
  - 0 V
  - +100 V
  - +200 V
  - +300 V
  - +400 V
  - +500 V
  - +600 V
  - +700 V

- **Dynodes:** Within each stage, dynodes are represented as inclined surfaces. Electrons are shown bouncing off these, multiplying at each stage due to secondary emission, highlighted with black lines indicating their path.

- **Vacuum:** The tube is labeled as a vacuum, ensuring that electron movement is unimpeded.

- **Output to Counter:** At the bottom, the final stage leads to an output connected to a counter, indicating the measurement of the amplified electron signal.

Overall, the diagram illustrates the process of electron multiplication in a photomultiplier tube, highlighting the role of each component in increasing the number of electrons generated from an initial photon or particle interaction.

\#\# Dataset metadata

Quantum Phenomena; Multi-Formula Reasoning, Physical Model Grounding Reasoning

\paragraph{Recorded answer/context.}
D: D: \textbackslash{}(10\textasciicircum{}8 \textbackslash{}text\{ eV\}\textbackslash{})

\paragraph{Figure/image path.}
/root/proposal\_for\_physic/science-mango/phyx\_mini\_run/phyx\_data/test\_image/539.png

\paragraph{Formalization target.}
create a compiling Lean file with sorry bodies at `PhyXMiniProblems/problem\_phyx\_mini\_0539.lean`. The Lean declarations must preserve the physical quantities, dimensions or dimensional roles, figure labels, governing-law hypotheses, and final relation expressed by this problem.
Use Mathlib/Physlib names found through LeanExplore where available. If a physics API is missing, introduce faithful local abstractions rather than scalar placeholder aliases.

\begin{theorem}[Physics formalization target]
\label{thm:physics:phyx_mini_0539:target}
\lean{PhyXMiniProblems.ProblemPhyXMini0539.problem_phyx_mini_0539}
\uses{def:physics:phyx-mini-0539:phyxminiproblems-problemphyxmini0539-energyinelectronvolts, def:physics:phyx-mini-0539:phyxminiproblems-problemphyxmini0539-sevendynodephotomultipliersetup, def:physics:phyx-mini-0539:phyxminiproblems-problemphyxmini0539-matchesphotomultiplierproblemdata, def:physics:phyx-mini-0539:phyxminiproblems-problemphyxmini0539-matchessuppliedphotomultiplierfigure, def:physics:phyx-mini-0539:phyxminiproblems-problemphyxmini0539-hasphysicalphotomultiplierparameters, def:physics:phyx-mini-0539:phyxminiproblems-problemphyxmini0539-satisfiesidealphotomultiplierlaws}
The assigned autoformalize agent should translate this physics problem into Lean declarations in the covered file, with theorem and lemma proof bodies written as `by sorry`.
\end{theorem}
\begin{proof}
This is an autoformalization task, not a proof task. Produce statements that can later be proved by the physics prover without weakening the physical model.
\end{proof}

% --- Archon declaration topology begin ---
\section{Declaration topology}
\begin{definition}[electric Potential Dimension]
  \label{def:physics:phyx-mini-0539:phyxminiproblems-problemphyxmini0539-electricpotentialdimension}
  \lean{PhyXMiniProblems.ProblemPhyXMini0539.electricPotentialDimension}
  Electric potential has dimension energy divided by electric charge
\end{definition}
\begin{proof}
  This is the declared model definition of electric Potential Dimension.
\end{proof}
\begin{definition}[Electric Potential Quantity]
  \label{def:physics:phyx-mini-0539:phyxminiproblems-problemphyxmini0539-electricpotentialquantity}
  \lean{PhyXMiniProblems.ProblemPhyXMini0539.ElectricPotentialQuantity}
  \uses{def:physics:phyx-mini-0539:phyxminiproblems-problemphyxmini0539-electricpotentialdimension}
  A signed, unit-independent physical electric potential
\end{definition}
\begin{proof}
  This is the declared model definition of Electric Potential Quantity.
\end{proof}
\begin{definition}[energy In Joules]
  \label{def:physics:phyx-mini-0539:phyxminiproblems-problemphyxmini0539-energyinjoules}
  \lean{PhyXMiniProblems.ProblemPhyXMini0539.energyInJoules}
  Read a physical energy in coherent-SI joules
\end{definition}
\begin{proof}
  This is the declared model definition of energy In Joules.
\end{proof}
\begin{definition}[energy In Electron Volts]
  \label{def:physics:phyx-mini-0539:phyxminiproblems-problemphyxmini0539-energyinelectronvolts}
  \lean{PhyXMiniProblems.ProblemPhyXMini0539.energyInElectronVolts}
  \uses{def:physics:phyx-mini-0539:phyxminiproblems-problemphyxmini0539-energyinjoules}
  Read a physical energy in electron volts
\end{definition}
\begin{proof}
  This is the declared model definition of energy In Electron Volts.
\end{proof}
\begin{definition}[electric Potential In Volts]
  \label{def:physics:phyx-mini-0539:phyxminiproblems-problemphyxmini0539-electricpotentialinvolts}
  \lean{PhyXMiniProblems.ProblemPhyXMini0539.electricPotentialInVolts}
  \uses{def:physics:phyx-mini-0539:phyxminiproblems-problemphyxmini0539-electricpotentialquantity}
  Read a physical electric potential in coherent-SI volts
\end{definition}
\begin{proof}
  This is the declared model definition of electric Potential In Volts.
\end{proof}
\begin{definition}[Tube Interior Medium]
  \label{def:physics:phyx-mini-0539:phyxminiproblems-problemphyxmini0539-tubeinteriormedium}
  \lean{PhyXMiniProblems.ProblemPhyXMini0539.TubeInteriorMedium}
  The material filling the photomultiplier enclosure
\end{definition}
\begin{proof}
  This is the declared model definition of Tube Interior Medium.
\end{proof}
\begin{definition}[Photomultiplier Figure]
  \label{def:physics:phyx-mini-0539:phyxminiproblems-problemphyxmini0539-photomultiplierfigure}
  \lean{PhyXMiniProblems.ProblemPhyXMini0539.PhotomultiplierFigure}
  Literal information carried by image 539. Electrode position 0 is the photocathode, and positions 1, ..., 7 are the seven dynodes in increasing potential order. The potential labels are scalar volt readouts
\end{definition}
\begin{proof}
  This is the declared model definition of Photomultiplier Figure.
\end{proof}
\begin{definition}[Seven Dynode Photomultiplier Setup]
  \label{def:physics:phyx-mini-0539:phyxminiproblems-problemphyxmini0539-sevendynodephotomultipliersetup}
  \lean{PhyXMiniProblems.ProblemPhyXMini0539.SevenDynodePhotomultiplierSetup}
  \uses{def:physics:phyx-mini-0539:phyxminiproblems-problemphyxmini0539-electricpotentialquantity, def:physics:phyx-mini-0539:phyxminiproblems-problemphyxmini0539-tubeinteriormedium, def:physics:phyx-mini-0539:phyxminiproblems-problemphyxmini0539-photomultiplierfigure}
  Independent physical quantities in the idealized tube. Dynode index 0 is the +100 V dynode, and index 6 is the last, +700 V, dynode. Electron counts and secondary yields are discrete; energy and potential retain their physical dimensions
\end{definition}
\begin{proof}
  This is the declared model definition of Seven Dynode Photomultiplier Setup.
\end{proof}
\begin{definition}[Matches Photomultiplier Problem Data]
  \label{def:physics:phyx-mini-0539:phyxminiproblems-problemphyxmini0539-matchesphotomultiplierproblemdata}
  \lean{PhyXMiniProblems.ProblemPhyXMini0539.MatchesPhotomultiplierProblemData}
  \uses{def:physics:phyx-mini-0539:phyxminiproblems-problemphyxmini0539-energyinelectronvolts, def:physics:phyx-mini-0539:phyxminiproblems-problemphyxmini0539-sevendynodephotomultipliersetup}
  Problem data stated in the prose: one initial photoelectron, 10.0 eV average liberation energy, a vacuum interior, and 100\% efficiency. The first-arrival field is an initial condition, not a conclusion about later multiplication
\end{definition}
\begin{proof}
  This is the declared model definition of Matches Photomultiplier Problem Data.
\end{proof}
\begin{definition}[Matches Supplied Photomultiplier Figure]
  \label{def:physics:phyx-mini-0539:phyxminiproblems-problemphyxmini0539-matchessuppliedphotomultiplierfigure}
  \lean{PhyXMiniProblems.ProblemPhyXMini0539.MatchesSuppliedPhotomultiplierFigure}
  \uses{def:physics:phyx-mini-0539:phyxminiproblems-problemphyxmini0539-electricpotentialinvolts, def:physics:phyx-mini-0539:phyxminiproblems-problemphyxmini0539-sevendynodephotomultipliersetup}
  All unambiguous data read from the supplied image: voltage labels 0, 100, ..., 700 V, their calibration to the physical electrodes, the scintillator and photocathode, seven dynode surfaces, multiplication tracks, the vacuum label, and the counter output
\end{definition}
\begin{proof}
  This is the declared model definition of Matches Supplied Photomultiplier Figure.
\end{proof}
\begin{definition}[Has Physical Photomultiplier Parameters]
  \label{def:physics:phyx-mini-0539:phyxminiproblems-problemphyxmini0539-hasphysicalphotomultiplierparameters}
  \lean{PhyXMiniProblems.ProblemPhyXMini0539.HasPhysicalPhotomultiplierParameters}
  \uses{def:physics:phyx-mini-0539:phyxminiproblems-problemphyxmini0539-energyinelectronvolts, def:physics:phyx-mini-0539:phyxminiproblems-problemphyxmini0539-electricpotentialinvolts, def:physics:phyx-mini-0539:phyxminiproblems-problemphyxmini0539-sevendynodephotomultipliersetup}
  Physical sign and range conditions, independent of the requested answer
\end{definition}
\begin{proof}
  This is the declared model definition of Has Physical Photomultiplier Parameters.
\end{proof}
\begin{definition}[Satisfies Ideal Photomultiplier Laws]
  \label{def:physics:phyx-mini-0539:phyxminiproblems-problemphyxmini0539-satisfiesidealphotomultiplierlaws}
  \lean{PhyXMiniProblems.ProblemPhyXMini0539.SatisfiesIdealPhotomultiplierLaws}
  \uses{def:physics:phyx-mini-0539:phyxminiproblems-problemphyxmini0539-energyinelectronvolts, def:physics:phyx-mini-0539:phyxminiproblems-problemphyxmini0539-electricpotentialinvolts, def:physics:phyx-mini-0539:phyxminiproblems-problemphyxmini0539-sevendynodephotomultipliersetup}
  The idealized laws used by the calculation are uniform across the seven dynodes: one electron accelerated through ΔV volts gains numerically ΔV electron volts of impact energy; the efficient part of that impact energy is the energy budget for freeing an integral number of secondary electrons; all emitted electrons propagate to the next dynode; and the counter receives the aggregate impact energy of the electrons arriving at the last dynode.
\end{definition}
\begin{proof}
  This is the declared model definition of Satisfies Ideal Photomultiplier Laws.
\end{proof}
\begin{lemma}[secondary Electron Yield At Every Dynode eq ten]
  \label{lem:physics:phyx-mini-0539:phyxminiproblems-problemphyxmini0539-secondaryelectronyieldateverydynode-eq-ten}
  \lean{PhyXMiniProblems.ProblemPhyXMini0539.secondaryElectronYieldAtEveryDynode_eq_ten}
  \uses{def:physics:phyx-mini-0539:phyxminiproblems-problemphyxmini0539-sevendynodephotomultipliersetup, def:physics:phyx-mini-0539:phyxminiproblems-problemphyxmini0539-matchesphotomultiplierproblemdata, def:physics:phyx-mini-0539:phyxminiproblems-problemphyxmini0539-matchessuppliedphotomultiplierfigure, def:physics:phyx-mini-0539:phyxminiproblems-problemphyxmini0539-hasphysicalphotomultiplierparameters, def:physics:phyx-mini-0539:phyxminiproblems-problemphyxmini0539-satisfiesidealphotomultiplierlaws}
  Every adjacent electrode gap is 100 V; perfect conversion of a 100 eV impact with a 10 eV liberation cost produces ten secondary electrons per incident electron at each dynode
\end{lemma}
\begin{proof}
  The conclusion follows from the governing relations and figure readouts named in the statement of secondary Electron Yield At Every Dynode eq ten.
\end{proof}
\begin{lemma}[arriving Electron Count At Last Dynode eq ten pow six]
  \label{lem:physics:phyx-mini-0539:phyxminiproblems-problemphyxmini0539-arrivingelectroncountatlastdynode-eq-ten-pow-six}
  \lean{PhyXMiniProblems.ProblemPhyXMini0539.arrivingElectronCountAtLastDynode_eq_ten_pow_six}
  \uses{def:physics:phyx-mini-0539:phyxminiproblems-problemphyxmini0539-sevendynodephotomultipliersetup, def:physics:phyx-mini-0539:phyxminiproblems-problemphyxmini0539-matchesphotomultiplierproblemdata, def:physics:phyx-mini-0539:phyxminiproblems-problemphyxmini0539-matchessuppliedphotomultiplierfigure, def:physics:phyx-mini-0539:phyxminiproblems-problemphyxmini0539-hasphysicalphotomultiplierparameters, def:physics:phyx-mini-0539:phyxminiproblems-problemphyxmini0539-satisfiesidealphotomultiplierlaws}
  Starting from one electron at the first dynode and multiplying by ten at the first six dynodes, 10\textasciicircum{}6 electrons arrive at the seventh dynode
\end{lemma}
\begin{proof}
  The conclusion follows from the governing relations and figure readouts named in the statement of arriving Electron Count At Last Dynode eq ten pow six.
\end{proof}
\begin{definition}[Answer Choice]
  \label{def:physics:phyx-mini-0539:phyxminiproblems-problemphyxmini0539-answerchoice}
  \lean{PhyXMiniProblems.ProblemPhyXMini0539.AnswerChoice}
  Labels printed beside the four energy choices
\end{definition}
\begin{proof}
  This is the declared model definition of Answer Choice.
\end{proof}
\begin{definition}[energy In Electron Volts]
  \label{def:physics:phyx-mini-0539:phyxminiproblems-problemphyxmini0539-answerchoice-energyinelectronvolts}
  \lean{PhyXMiniProblems.ProblemPhyXMini0539.AnswerChoice.energyInElectronVolts}
  \uses{def:physics:phyx-mini-0539:phyxminiproblems-problemphyxmini0539-energyinelectronvolts, def:physics:phyx-mini-0539:phyxminiproblems-problemphyxmini0539-answerchoice}
  Displayed scalar energy for each choice, in electron volts
\end{definition}
\begin{proof}
  This is the declared model definition of energy In Electron Volts.
\end{proof}
\begin{definition}[recorded Dataset Answer]
  \label{def:physics:phyx-mini-0539:phyxminiproblems-problemphyxmini0539-recordeddatasetanswer}
  \lean{PhyXMiniProblems.ProblemPhyXMini0539.recordedDatasetAnswer}
  \uses{def:physics:phyx-mini-0539:phyxminiproblems-problemphyxmini0539-answerchoice}
  The answer label recorded by the source dataset
\end{definition}
\begin{proof}
  This is the declared model definition of recorded Dataset Answer.
\end{proof}
% --- Archon declaration topology end ---
% --- Archon physics formalization source end ---
